% Drop-in main text; Figure 4 is supplied separately after its revision.
\subsection{Why joint training acquires useful feature scales slowly}
\label{sec:feature-acquisition}

Joint training must acquire useful slopes and centers together.
After five million updates on mixed sine at $W=512$, joint Adam reaches
median relative training error $1.81\times10^{-3}$, versus
$6.62\times10^{-7}$ with supplied $\lambda=1/4$ features;
joint GD remains at $0.244$ (Figure~\ref{fig:note-training}b).
Numerically resolved readout refits barely improve the learned Adam
dictionary, and increasing its slopes at fixed centers does not close the
gap (Appendix~\ref{app:feature-evidence}).

After coarse fitting, feature acquisition depends on nonlinear output error.
The cubic term of tanh couples to the target only through its projection
onto polynomials of degree at most three; weak projection suppresses this
leading target contribution. Bounding that coupling, the network's own
nonlinear output, and higher-order remainders controls the growth of slopes,
biases, and readouts together. For populations with small total parameter
norm and controlled concentration, this yields an output-error lower bound.

\begin{theorem}[Output error under slow feature acquisition]
\label{thm:feature-output}
Let $f\ne0$ be the sampled target and $f_\perp$ its component orthogonal
to affine functions. Under the population-concentration, target-coupling,
and coarse-tracking bounds of Appendix~\ref{app:feature-theory},
joint GD satisfies
\begin{equation}
 \frac{\|f-q_n\|_{2,Z}}{\|f\|_{2,Z}}
 \ge
 \frac{[\|f_\perp\|_{2,Z}-\mathcal Q_n(B_n)]_+}{\|f\|_{2,Z}}.
 \label{eq:feature-main-error}
\end{equation}
Here $B_n$ is an explicit scalar upper bound on the joint norm of hidden
slopes, biases, and readouts, and $\mathcal Q_n(B_n)$ bounds their
non-affine output. The same comparison gives
$\overline\lambda_n\le hB_n/\sqrt W$, with reference spacing $h$ and
$\overline\lambda_n=W^{-1}\sum_j h|\gamma_{j,n}|$.
The recurrence, coefficients, and precise conditions are given in the appendix.
\end{theorem}

\paragraph{Proof sketch.}
Separate the gradient preserving affine output from the remaining tracking
gradient. Global tanh remainder bounds express nonlinear output and target
coupling through population moments. An upper bound on parameter movement
propagates $B_n$ by a scalar recurrence; subtracting the resulting
non-affine output bound from $\|f_\perp\|_{2,Z}$ yields
\eqref{eq:feature-main-error}. All parameters evolve, and the comparison
allows positive growth.

On ten post-transient GD continuations across six target families, minimum
evaluated error floors range from $0.427$ to $0.912$ over 100,000 further updates.
These checks use structural quantities sampled along training
(Appendix~\ref{app:feature-evidence}); they establish a useful conditional
mechanism in that regime. Figure~\ref{fig:note-training}'s longer Adam/GD
runs establish the broader empirical precision gap. The construction
supplies the geometry and readout directly, avoiding the need to acquire
them through this slow process.
